\begin{figure*}[t]
\centering
\begin{tcolorbox}[title=Cheat Sheet for Disambiguation QA (1/2), colback=gray!10, colframe=black!70, boxrule=0.5pt, top=0pt, bottom=0pt, left=0pt, right=0pt, arc=0.8mm, width=1.0\textwidth]
\scriptsize
\begin{lstlisting}
---

# Pronoun Antecedent Cheat Sheet (for Difficult Cases)

## 1. **General Reasoning Steps**
- **Identify all possible antecedents** for the pronoun.
- **Substitute each antecedent** into the sentence to see if it makes sense.
- **Consider the context and logic** of the sentence: Who is likely to perform the action or possess the attribute?
- **Check for grammatical cues**: number (singular/plural), gender, and role in the sentence.
- **If both options are equally plausible and the sentence gives no extra clues, mark as ambiguous.**

---

## 2. **Common Patterns and How to Resolve Them**

### A. **"X told Y that [pronoun]..."**
- Usually, the pronoun refers to **Y** if the information is about Y (e.g., advice, diagnosis, payment).
- If the information is about X's own actions or status, it refers to **X**.
- **Tip:** Would it make sense for X to inform Y about Y's own actions? Usually not, unless it's advice or a warning.

### B. **"X did something to Y because [pronoun]..."**
- The pronoun can refer to either X or Y.
- **Test both:** Substitute both and see which makes more logical sense.
- If both are plausible, **mark as ambiguous**.

### C. **"X and Y discuss [pronoun]'s Z"**
- If both X and Y could logically possess Z, and the sentence gives no further context, **mark as ambiguous**.
- If only one is likely to possess Z (e.g., "culinary training" is more likely the chef's), pick that one.

### D. **"X called Y and asked [pronoun] to do Z"**
- The pronoun usually refers to **Y** (the person being asked to do something).
- If it would be odd for X to ask themselves, it's almost always Y.

### E. **"X met with Y at [pronoun]'s office"**
- If both X and Y could be the owner of the office, and the sentence gives no clue, **mark as ambiguous**.
- If only one is plausible (e.g., meeting a director at the director's office), pick that one.

### F. **"X did something with Y because [pronoun] [verb/attribute]"**
- If the verb/attribute fits both X and Y, and both are plausible, **mark as ambiguous**.
- If only one makes sense (e.g., "focuses on code" fits developer, not writer), pick that one.

### G. **Possessive Constructions ("the writer and [pronoun] friends")**
- The possessive pronoun almost always refers to the first noun ("the writer and her friends" = the writer's friends).
- If the pronoun could refer to more than one noun, but only one makes sense, pick that one.

---
\end{lstlisting}
\end{tcolorbox}
\caption{
An example of cheat sheet generated for Disambiguation QA.
This is the first half of the cheat sheet.
}
\label{fig:cheat_sheet_disamb_1}
\end{figure*}